% #############################################################################
% This is Chapter 6
% !TEX root = ../main.tex
% #############################################################################
% Change the Name of the Chapter i the following line
\fancychapter{Conclusion}
% The following line allows to ref this chapter
\label{chap:conclusion}

% NOTE: skeleton + bullet-point to-be content only (no prose yet).

% #############################################################################
\section{Summary of Contributions}\label{sec:concl-summary}
\begin{itemize}
    \item Recap the five contributions from \Cref{sec:intro-contributions} in completed tense, each cross-referencing the chapter/table that substantiates it.
\end{itemize}

% #############################################################################
\section{Answers to Research Questions}\label{sec:concl-rqs}
\begin{itemize}
    \item One paragraph per RQ (\Cref{sec:eval-rqs}), each closing with the evidence (table/figure ref) that answers it.
\end{itemize}

% #############################################################################
\section{System Limitations}\label{sec:concl-limitations}
\begin{itemize}
    \item Hard reflection-flood sources and IAT-only perturbations (real, measured limitation).
    \item Purification evaluated through an own-TCN oracle, not an independent NIDS.
    \item Any limitation specific to the adaptive component --- to be filled in later.
    \item Feature-space scope: nPrint + IPAT + Direction; no payload/DPI features (relevant given the TLS/QUIC encryption trend).
\end{itemize}

% #############################################################################
\section{Future Work}\label{sec:concl-future}
\begin{itemize}
    \item Independent-NIDS purification validation (closing the own-oracle limitation).
    \item Maturing/extending the adaptive (RL) component.
    \item Feature-space expansion / payload-aware representations for encrypted-traffic robustness.
    \item Real-time deployment study at line rate (100\,Gbps+) beyond the simulated latency benchmark.
\end{itemize}
